\documentclass[11pt]{article}

\usepackage[margin=1.1in]{geometry}
\usepackage{amsmath}
\usepackage{amssymb}
\usepackage{amsthm}
\usepackage{booktabs}
\usepackage{microtype}
\usepackage[round,authoryear]{natbib}
\usepackage[colorlinks=true,allcolors=blue]{hyperref}

\makeatletter
\renewcommand\paragraph{\@startsection{paragraph}{4}{\z@}%
  {3.25ex \@plus1ex \@minus.2ex}{-1em}%
  {\normalfont\normalsize\itshape}}
\makeatother

\newtheoremstyle{itheads}{\topsep}{\topsep}%
  {\itshape}{}{\itshape}{.}{ }{\thmname{#1}\thmnumber{ #2}%
  \thmnote{ (#3)}}
\theoremstyle{itheads}
\newtheorem{proposition}{Proposition}
\newtheorem{lemma}{Lemma}

\newcommand{\E}{\mathbb{E}}

\title{\textbf{Wavefront Pruning in Budgeted Brownian Races}}
\author{Peter Cotton}
\date{\today}

\begin{document}
\maketitle

\begin{abstract}
\noindent In a budgeted Brownian race, a controller watches a cloud of
Brownian paths, pays per unit time for every path kept alive, may
irreversibly prune, and spends a finite path-time budget to maximize the
expected terminal maximum. We solve the mean-field version: a Poissonized
population under per-path controls, with the budget enforced in
expectation. The expected maximum has an exact layer-cake form whose
gradient is a pivotal value, and for a fixed price on path-time the
linearized best response is an obstacle problem whose keep sets are upper
tails: the \emph{wavefront policy}, which retains a path while its
propagated future pivotal value exceeds its carrying cost. A fully
corrective Frank--Wolfe method over wavefront policies, with dual
bisection on the price, solves the budget problem and certifies its
answer, bounding suboptimality by the Frank--Wolfe gap plus a
complementary slackness term. In the committed example the adaptive
policy improves the expected maximum by 19\% over the best single
screening date and by 30\% over static thinning at the same expected
budget; the finite-$n$ race with a pathwise budget is left to future
work.
\end{abstract}

\section{Introduction}

A search that runs many exploratory processes in parallel faces one
recurring decision: which of them to keep paying for. The processes drift,
their ranking changes, killing is irreversible, and only the best survivor
is paid at the end. This paper studies the cleanest continuous model of
that decision, which we call a \emph{budgeted Brownian race}: a controller
observes Brownian paths, pays per unit time for each surviving path, may
irreversibly prune, and allocates a finite path-time budget to maximize
the expected terminal maximum over a deterministic fallback.

The distinctive object is the \emph{wavefront}: the narrow region near the
future head of the race where one more path can still be terminally
pivotal. The optimal policy has a one-line description: retain a path
while its propagated future pivotal value exceeds its carrying cost. The
body of the paper makes the description exact in the mean-field limit and
computable at scale.

The name is chosen against three neighbors. Branching Brownian motion with
selection retains a prescribed number of rightmost particles; here nothing
branches, the survivor count is endogenous, and the resource is cumulative
path-time rather than a headcount. Finite-fuel stochastic control is the
classification home of the constraint, not a description of the
terminal-maximum payoff. Beam search fixes its width in advance; wavefront
pruning lets the width float under the budget.

\section{The model}

Optional paths form a Poisson cloud with initial intensity measure
$\mu_0$, and each alive path follows an independent Brownian motion. At
each decision time a path is either killed irreversibly or kept for one
more interval; keeping one path for $dt$ costs $dt$. The budget $B$
constrains expected aggregate path-time. The terminal payoff is the
maximum of a deterministic fallback $f$ and the surviving paths.

Controls are per-path and population-blind: each path's keep-or-kill
decision depends on its own state and time, plus independent
randomization, never on the realized population. Under this restriction
the surviving cloud remains Poisson at every date, by independent
thinning, and the terminal intensity $m$ is the object of choice.

The restriction has a price a one-step example makes exact. Start
$\mathrm{Poisson}(100)$ paths at zero with budget one and fallback zero.
Independent thinning to intensity one pays about $0.3469$; a controller
who sees the population and keeps exactly one path whenever any exist
pays $(1-e^{-100})/\sqrt{2\pi} \approx 0.3989$ on the same expected budget. The
values in this paper are optimal within the population-blind class.

\section{The value of the terminal cloud}

Fix a terminal grid $x_1 < \dots < x_J$ and terminal intensity $m$ with
upper-tail intensity $R_j = \sum_{k \ge j} m_k$. The expected maximum is
exactly
\begin{equation}
J(m) \;=\; f + \sum_{x_j > f} \Delta_j \left(1 - e^{-R_j}\right),
\label{eq:layercake}
\end{equation}
where $\Delta_j$ is the gap to the preceding attainable level. In
English: each layer above the fallback pays its width times the
probability that at least one survivor clears it. $J$ is concave in $m$,
and its gradient is the \emph{pivotal value}
\begin{equation}
\frac{\partial J}{\partial m_j} \;=\; \E\big[(x_j - M)^+\big],
\label{eq:pivotal}
\end{equation}
the expected improvement from placing one more infinitesimal path at
$x_j$, where $M$ is the current maximum.

\section{Wavefront policies}

For a shadow price $\lambda \ge 0$ on path-time, linearizing $J$ at the
current intensity gives a one-particle obstacle problem,
\begin{equation}
v_N = \E[(x - M)^+], \qquad
v_k(x) = \max\!\big(0,\; \E[\,v_{k+1}(X_{k+1}) \mid X_k = x\,] -
\lambda\, dt\big),
\label{eq:obstacle}
\end{equation}
whose keep set at each step is an upper tail of the state space: one
kill-below cutoff per time step. We call this the wavefront policy. The
optimal population policy is in general a mixture of a few wavefront
atoms, and the mixture is implementable per path: each path independently
follows one atom, and by Poisson thinning the population intensity is
exactly the mixture's.

\section{The algorithm and its certificate}

A fully corrective Frank--Wolfe (column-generation) method repeatedly
computes the pivotal-value gradient \eqref{eq:pivotal}, solves the
obstacle problem \eqref{eq:obstacle} for the best-response wavefront
policy, and re-optimizes the mixture of active atoms; bisection in
$\lambda$ enforces the budget. The scheme is a discrete counterpart of
the coupled obstacle/Fokker--Planck KKT system.

The Frank--Wolfe gap alone does not certify accuracy at the requested
budget, because the returned policy may underspend. The correct bound
carries a complementary slackness term.

\begin{lemma}[Budget suboptimality certificate]
Let $C(m)$ denote expected path-time, $V(B) = \max\{J(m) : C(m) \le B\}$,
and let $m$ be any feasible intensity, with $g$ the Frank--Wolfe gap of
the Lagrangian $J - \lambda C$ at $m$. Then for every $\lambda \ge 0$,
\[
V(B) - J(m) \;\le\; g + \lambda\,(B - C(m)).
\]
\end{lemma}

\begin{proof}
By concavity of $J$, the Lagrangian optimum satisfies $L^\star \le
L(m) + g$. For any feasible $m'$, $J(m') \le L(m') + \lambda B \le
L^\star + \lambda B \le L(m) + g + \lambda B$, and $L(m) = J(m) -
\lambda C(m)$.
\end{proof}

The implementation returns a policy with $C(m) \le B$, thinning with the
kill-all policy when the bisection overspends, and reports the two terms
of the certificate separately. The slackness term is not small in
practice: in one representative solve it exceeds the gap by a factor of
seventy.

\section{Related work}

\paragraph{Selection in branching systems.}
\citet{brunet1997} quantified what truncating a population does to its
leading edge; \citet{berestycki2013} worked out absorption at a moving
barrier as a selection mechanism; \citet{maillard2016} pinned the speed
and fluctuations of $N$-particle branching Brownian motion under spatial
selection. In all three the population size is exogenous. The budgeted
race replaces the fixed-$N$ constraint with cumulative path-time, and
nothing branches.

\paragraph{Budgets as fuel.}
\citet{benes1980} solved the original finite-fuel control problems in
closed form. The label classifies the constraint here, a bounded
cumulative control effort, but not the payoff.

\paragraph{Search, indices, and optionality.}
\citet{weitzman1979} opened boxes in reservation-value order;
\citet{gittins1979} priced the next unit of attention with an index;
\citet{audibert2010} studied pure exploration under a hard budget of
pulls. The bandit literature closest in payoff is the max-$k$-armed
bandit of \citet{cicirello2005} and the extreme bandits of
\citet{carpentier2014}, where the criterion is the maximum reward found
rather than the sum. The race is their continuous cousin: arms drift as
diffusions, attention is metered in path-time, and killing is
irreversible.

\paragraph{Mean-field control.}
\citet{lasry2007} coupled a value function to a density through a
consistency condition; \citet{carmona2018} treat control of
McKean--Vlasov dynamics, the honest label for the population problem
here. The race's version replaces the HJB by an obstacle problem, since
the only control is stopping.

\paragraph{The algorithmic tools.}
\citet{frank1956} linearize and call an oracle; \citet{jaggi2013}
supplies the duality-gap toolkit. The oracle here is the obstacle
problem, and the atoms it returns are wavefront policies.

\paragraph{Budgeted rollouts of language models.}
Best-of-$n$ sampling is a budgeted race: partial generations drift,
compute is metered per token kept alive, and only the best completion is
paid \citep{chen2021}. \citet{sun2024} kill low-reward generations
mid-decode and reach best-of-$n$ quality at a fraction of the cost;
\citet{bi2026} learn a prefix-level pruning signal for parallel
reasoning under a fixed budget; \citet{kalayci2025} stop sequential
generation by Pandora's-box reasoning; \citet{nomand2026} allocate a
rollout budget by a two-threshold stopping rule. Each learns, scores,
or proves its kill rule in a one-shot or sequential relaxation. None
models the drifting value of a partial trajectory as a diffusion and
derives the kill boundary; that is the gap this paper fills, and the
learned rules are the ones its wavefront would predict.
\citet{yao2023} and \citet{snell2024} frame the wider test-time-compute
allocation problem; \citet{malladi2026} train for coverage, shaping the
distribution whose maximum is paid. The estimation side, predicting the
expected maximum of a batch from few samples, is treated in a companion
note \citep{cotton2026passk}.

\section{Numerical certificates}

Every numerical claim is backed by a committed, seeded script in the
repository\footnote{\texttt{github.com/microprediction/brownianbandit};
live demos at \texttt{brownianbandit.microprediction.org}.}. In the
committed example (horizon $1$, initial intensity $100$, budget $10$,
fallback $0$, $100$ time steps, $601$ state points) the adaptive policy
reaches an expected maximum of $1.9600$, at a cost of exactly the
budget and with a certificate below $2 \times 10^{-6}$, against
$1.6429$ for the best one-shot screening policy and $1.5086$ for static
random thinning. A $100{,}000$-trial Monte Carlo puts the adaptive
policy at $1.9587$ with standard error $0.0021$. The analytic gradient of \eqref{eq:layercake}
matches finite differences to $3 \times 10^{-7}$; spent path-time is
monotone in the shadow price; refining the grid moves the objective by
less than $0.08$. A dependency-free JavaScript port runs the same solve
in the browser and is held to the Python at stated tolerances by a
committed parity test.

\section{Closing}

The mean-field, population-blind race is solved, certified, and cheap.
Three problems remain open, in increasing order of difficulty: the value
of population awareness beyond the one-step example of Section~2; the
finite-$n$ race under a pathwise hard budget, where the Poisson structure
is lost; and the continuous-time limit, where we expect the wavefront
cutoff to become a free boundary and the pivotal-value flux across it a
collision local time. We leave to future work the correlated race, where
paths share factors and the pruning decision must price a rival's
strength.

\begin{thebibliography}{99}

\bibitem[Audibert et al.(2010)]{audibert2010}
Audibert, J.-Y., Bubeck, S., and Munos, R. (2010).
Best arm identification in multi-armed bandits.
In \emph{COLT}.

\bibitem[Bene\v{s} et al.(1980)]{benes1980}
Bene\v{s}, V. E., Shepp, L. A., and Witsenhausen, H. S. (1980).
Some solvable stochastic control problems.
\emph{Stochastics}, 4(1):39--83.

\bibitem[Berestycki et al.(2013)]{berestycki2013}
Berestycki, J., Berestycki, N., and Schweinsberg, J. (2013).
The genealogy of branching Brownian motion with absorption.
\emph{Annals of Probability}, 41(2):527--618.

\bibitem[Bi et al.(2026)]{bi2026}
Bi, J., Luo, T., Du, W., Tang, Z., and Wang, B. (2026).
Cut your losses! Learning to prune paths early for efficient parallel
reasoning.
arXiv:2604.16029.

\bibitem[Brunet and Derrida(1997)]{brunet1997}
Brunet, E., and Derrida, B. (1997).
Shift in the velocity of a front due to a cutoff.
\emph{Physical Review E}, 56(3):2597--2604.

\bibitem[Carmona and Delarue(2018)]{carmona2018}
Carmona, R., and Delarue, F. (2018).
\emph{Probabilistic Theory of Mean Field Games with Applications},
vols.\ I--II. Springer.

\bibitem[Carpentier and Valko(2014)]{carpentier2014}
Carpentier, A., and Valko, M. (2014).
Extreme bandits.
In \emph{NeurIPS}.

\bibitem[Chen et al.(2021)]{chen2021}
Chen, M., et al. (2021).
Evaluating large language models trained on code.
arXiv:2107.03374.

\bibitem[Cicirello and Smith(2005)]{cicirello2005}
Cicirello, V. A., and Smith, S. F. (2005).
The max $k$-armed bandit: a new model of exploration applied to search
heuristic selection.
In \emph{AAAI}.

\bibitem[Cotton(2026)]{cotton2026passk}
Cotton, P. (2026).
Posterior-predictive pass-at-$k$.
Working note, \texttt{brownianbandit.microprediction.org}.

\bibitem[Frank and Wolfe(1956)]{frank1956}
Frank, M., and Wolfe, P. (1956).
An algorithm for quadratic programming.
\emph{Naval Research Logistics Quarterly}, 3:95--110.

\bibitem[Gittins(1979)]{gittins1979}
Gittins, J. C. (1979).
Bandit processes and dynamic allocation indices.
\emph{Journal of the Royal Statistical Society, Series B},
41(2):148--177.

\bibitem[Jaggi(2013)]{jaggi2013}
Jaggi, M. (2013).
Revisiting Frank--Wolfe: projection-free sparse convex optimization.
In \emph{ICML}.

\bibitem[Kalayci et al.(2025)]{kalayci2025}
Kalayci, Y., Raman, V., and Dughmi, S. (2025).
Optimal stopping vs best-of-$N$ for inference time optimization.
arXiv:2510.01394.

\bibitem[Lasry and Lions(2007)]{lasry2007}
Lasry, J.-M., and Lions, P.-L. (2007).
Mean field games.
\emph{Japanese Journal of Mathematics}, 2(1):229--260.

\bibitem[Maillard(2016)]{maillard2016}
Maillard, P. (2016).
Speed and fluctuations of $N$-particle branching Brownian motion with
spatial selection.
\emph{Probability Theory and Related Fields}, 166:1061--1173.

\bibitem[Malladi et al.(2026)]{malladi2026}
Malladi, S., Jelassi, S., Foster, D., Ash, J. T., and Krishnamurthy, A.
(2026).
TailSFT: filtered fine-tuning improves post-training performance.
arXiv:2608.25756.

\bibitem[Nomand et al.(2026)]{nomand2026}
Nomand, P., Voss, E., Hale, M., and Reyes, S. (2026).
Early verdicts, better budgets: sequential adaptive rollout allocation
for compute-efficient RLVR.
arXiv:2607.26253.

\bibitem[Snell et al.(2024)]{snell2024}
Snell, C., Lee, J., Xu, K., and Kumar, A. (2024).
Scaling LLM test-time compute optimally can be more effective than
scaling model parameters.
arXiv:2408.03314.

\bibitem[Sun et al.(2024)]{sun2024}
Sun, H., Haider, M., Zhang, R., Yang, H., Qiu, J., Yin, M., Wang, M.,
Bartlett, P., and Zanette, A. (2024).
Fast best-of-$N$ decoding via speculative rejection.
In \emph{NeurIPS}.

\bibitem[Weitzman(1979)]{weitzman1979}
Weitzman, M. L. (1979).
Optimal search for the best alternative.
\emph{Econometrica}, 47(3):641--654.

\bibitem[Yao et al.(2023)]{yao2023}
Yao, S., Yu, D., Zhao, J., Shafran, I., Griffiths, T. L., Cao, Y., and
Narasimhan, K. (2023).
Tree of Thoughts: deliberate problem solving with large language models.
In \emph{NeurIPS}.

\end{thebibliography}

\end{document}
